\subsection{A further reduction to 663 bits: query native noncollapse}
\label{sec:e663}

The truth-value query $N(d,d)=d$ is not the shortest available ground
query with the required calibration.  Keep all seven equations of
$\mathcal E_{667}$ and the complete ordered signature unchanged, but
replace the designated query by
\begin{equation}
                            Q(d)=d.
 \label{eq:e663-query}
\end{equation}
The resulting complete consistency packet has \textbf{663 bits}.
This reduction uses the native coordinate equations, independently of
the Boolean truth-forcing optimization.

\paragraph{Query-equivalence theorem.}
For every algebra satisfying the seven equations, each of
$Q(d)=d$ and $N(d,d)=d$ holds if and only if its carrier is a
singleton.  Consequently the two queries have the same equational
derivability status from the seven equations.

\paragraph{Proof.}
Suppose $Q(d)=d$.  The native head and tail equations, with the
tail index $j=d$, give for every $x,s$
\[
 x=V(P(x,s),d)=V(P(x,s),Q(d))=V(s,d).
\]
Fixing $s=d$ shows that every element $x$ equals $V(d,d)$, so the
carrier is a singleton.  No Boolean identity or semantic quotient
is used in this implication.

Sh1 makes $N$ NAND in a Boolean algebra on this same carrier.
If $N(d,d)=d$, then $\neg d=d$, which implies
$0=d\wedge\neg d=d$ and $1=d\vee\neg d=d$.
Thus $0=1$ and the carrier is again a singleton.  Conversely both
equalities hold in every singleton algebra.  The two equalities
therefore hold in exactly the same models of the seven equations.
Equational soundness and completeness give equivalence of their
derivability.  Equivalently, use the countable ground-term congruence
model to witness either nonderivability; the constant $d$ guarantees
that this model is nonempty.
\hfill$\square$

The E667 consistency theorem therefore gives, with its same inherited
ordinary raw-EPSC/ZFC calibration,
\[
 \mathrm{ZFC}\vdash
 \left(\mathcal E_{667}\nvdash Q(d)=d
              \ \longleftrightarrow\ \operatorname{Con}(\mathrm{ZFC})\right).
\]
The label $\mathcal E_{663}$ denotes this new \emph{packet}; the
seven-equation axiom set itself is unchanged.  This distinction matters
because the four-bit saving is in the encoded goal, not the axioms.

The header still has 43 bits and the equations have 608 bits.  The new
endpoints have three term nodes, so they require 12 rather than 16 bits:
\[
                       43+608+12=\boxed{663\text{ bits}}.
\]
Packed transport occupies 83 bytes, with one trailing zero padding bit.
The exact bit-character SHA-256, excluding the final newline, is
\begin{center}\scriptsize\ttfamily
6c3b72c39c81aa5bbf028e654c8056e7e86376cdeecdd56943e5c3e23de1a215.
\end{center}
Run \texttt{python3 -B e663/build\_e663.py} from the E667 companion
directory.  It reconstructs the pinned E667 packet, verifies that the
header and all axiom bits are literally unchanged, and emits
\texttt{E663.bits}, \texttt{E663.bin}, \texttt{E663.bitlength},
\texttt{E663.json}, and an exact verification receipt.
All 663 strict truncations and three malformed controls are rejected
by the independent decoder.

\paragraph{A limited endpoint optimum.}
The signature has exactly one constant, $d$.  Thus its only one-node
ground term is $d$.  Two syntactically distinct ground endpoints must
have at least three nodes in total: one endpoint has at least two
nodes, and the other has at least one.  The pair $Q(d),d$ attains this
bound.  Since every node costs four bits in the unchanged codec,
\textbf{12 endpoint bits are optimal among distinct ground endpoint
pairs in this fixed signature and codec}.  This is not a claim that
663 total bits are globally minimal: the axioms, signature, codec,
or interpreter could be redesigned.  No machine-state bound is
inferred from the shorter query.
